% Libro 3 · Mundo 3 · Información explícita (app: pu-explicita)
\mundo{Información explícita}{pu-explicita}{La mitad de las preguntas de lectura se responden con algo que el texto
\emph{dice}. Parece fácil, y por eso se pierden: la respuesta correcta casi nunca copia las palabras del texto, sino que
las dice de otra manera, y las falsas copian palabras pero cambian el sentido. Este mundo entrena el ojo para la
precisión.}{%
\item \textbf{3.1} Ubicar datos
\item \textbf{3.2} Paráfrasis: decir lo mismo con otras palabras
\item \textbf{3.3} Lo que el texto dice y lo que no
\item \textbf{3.4} Quién dice qué}

% ------------------------------------------------------------------
\leccion{pu-ie-1}{Ubicar datos}{Juego de subrayar el trozo exacto que responde.}

\begin{reto}
«El embalse abastece a unas 400\,000 personas; en temporada seca baja hasta el 35\,\% de su capacidad, y por debajo del
30\,\% se declara alerta.» Según el fragmento, ¿cuándo se declara la alerta?
\opcionesl{Cuando el embalse llega al 35\,\%.}{Cuando baja de 30\,\%.}{En toda temporada seca.}{Cuando hay 400\,000 personas.}
\end{reto}

\tecnica{La técnica: palabra ancla y lectura fina}
Para ubicar un dato no se relee todo: se busca una \clave{palabra ancla} de la pregunta (un nombre, una fecha, un término
técnico) y se lee con lupa la oración donde aparece, y la anterior y la siguiente. Luego se compara con las opciones
\emph{palabra por palabra}: en las preguntas de datos, los distractores suelen cambiar un número, un «antes» por un
«después», un «algunos» por un «todos».

\figura{ancla}{La palabra ancla lleva directo a la oración que responde.}

\begin{regla}
\begin{itemize}
  \item Subraya en la pregunta la palabra ancla y búscala en el texto.
  \item Lee la oración completa y sus vecinas: el dato puede estar repartido.
  \item Compara cifras, tiempos y cuantificadores con lupa.
\end{itemize}
\end{regla}

\begin{texto}{El viaje del agua}{Artículo · texto original escrito para este libro (el municipio y sus cifras son imaginarios)}
\pa Abrir la llave en Villa Clara parece un gesto sencillo, pero detrás hay un recorrido de más de cuarenta kilómetros. El
agua nace en el páramo de La Ceja Alta, a 3\,400 metros de altura, donde los frailejones la retienen como esponjas y la
sueltan poco a poco hacia las quebradas.

\pa Esas quebradas alimentan el embalse de San Rafael, construido en 1987, que hoy abastece a unas 400\,000 personas. En
temporada de lluvias suele estar por encima del 80\,\% de su capacidad; en temporada seca puede bajar hasta el 35\,\%, y
cuando cae por debajo del 30\,\% la empresa de acueducto declara la alerta y pide ahorrar.

\pa Del embalse, el agua baja por gravedad hasta la planta de tratamiento, donde pasa por cuatro etapas: primero se le
agregan sustancias que hacen que la tierra y las partículas se agrupen; después esas partículas se dejan asentar en el
fondo de grandes tanques; luego el agua atraviesa filtros de arena y grava; y al final se desinfecta con cloro.

\pa Solo entonces entra a la red de tuberías. Los técnicos calculan que alrededor de una cuarta parte del agua tratada se
pierde en fugas antes de llegar a las casas, un problema que la empresa intenta reducir cambiando los tubos más viejos,
algunos de los cuales tienen más de sesenta años.
\end{texto}

\subsection{Preguntas comentadas}
\begin{comentada}{Según el texto, ¿en qué orden ocurren las etapas del tratamiento?}
\mal{A}{filtrar, asentar, agrupar, desinfectar}{\emph{Dice otra cosa}: altera el orden.}
\bien{B}{agrupar partículas, dejarlas asentar, filtrar, desinfectar}{Es el orden exacto del \S\,3 («primero…, después…,
luego…, al final…»).}
\mal{C}{desinfectar, filtrar, asentar, agrupar}{Invierte el orden.}
\mal{D}{asentar, agrupar, desinfectar, filtrar}{\emph{Dice otra cosa}.}
\end{comentada}

\begin{trampa}
Las cifras cercanas: 35\,\% y 30\,\%; 80\,\% y «una cuarta parte». Un distractor típico pone una cifra del texto en el
lugar de otra.
\end{trampa}

\taller
\piso{Como en el examen}{opción múltiple}
\begin{ejercicios}
\item Según el texto, el embalse fue construido en
\opcionesd{1987}{1960}{3\,400}{no se dice}
\sol{A. \S\,2.}
\item ¿Qué fracción del agua tratada se pierde en fugas, según los técnicos?
\opcionesd{la mitad}{alrededor de la cuarta parte}{el 30\,\%}{el 35\,\%}
\sol{B. \S\,4.}
\item ¿Qué papel cumplen los frailejones según el texto?
\opcionesl{Purifican el agua con cloro.}{Retienen el agua y la sueltan poco a poco.}{Producen la lluvia del páramo.}{Impiden las fugas.}
\sol{B. \S\,1.}
\item Algunos tubos de la red tienen
\opcionesd{menos de diez años}{cuarenta años}{más de sesenta años}{la edad del embalse}
\sol{C. \S\,4.}
\end{ejercicios}

\begin{retoresuelto}
El fragmento dice «por debajo del 30\,\% se declara alerta». El 35\,\% es otro dato (el nivel al que puede bajar).
\textbf{Opción B.}
\end{retoresuelto}
\bitacora

% ------------------------------------------------------------------
\leccion{pu-ie-2}{Paráfrasis}{Juego de emparejar una frase con su versión en otras palabras.}

\begin{reto}
«No todas las abejas pican.» ¿Cuál es una paráfrasis correcta?
\opcionesd{Ninguna abeja pica.}{Algunas abejas no pican.}{Las abejas no pican.}{Todas las abejas pican poco.}
\end{reto}

\tecnica{La técnica: mismo sentido, otras palabras}
Una \clave{paráfrasis} dice lo mismo que el original con otras palabras. Para que sea correcta debe conservar:
\begin{itemize}
  \item quién hace qué (los papeles no se intercambian);
  \item los \emph{cuantificadores} («algunos», «la mayoría», «todos», «ninguno»);
  \item el \emph{grado} de seguridad («podría», «es probable», «es seguro»);
  \item las relaciones (causa, condición, contraste).
\end{itemize}
Las opciones falsas cambian justamente uno de esos elementos, mientras conservan casi todas las palabras.

\figura{parafrasis}{Una paráfrasis conserva el sentido; el distractor cambia una pieza.}

\begin{regla}
\begin{itemize}
  \item Compara el sentido, no las palabras: la correcta suele usar sinónimos.
  \item Desconfía de la opción que copia el texto pero cambia «algunos» por «todos» o «puede» por «es».
  \item Pregúntate: si el texto fuera verdadero, ¿la opción tendría que serlo también?
\end{itemize}
\end{regla}

\begin{texto}{Abejas que no pican}{Divulgación · texto original escrito para este libro}
\pa Cuando se habla de abejas, casi todos piensan en la abeja de miel común, de origen europeo y africano, con su aguijón y
sus panales de cera en cajas de madera. Pero en los bosques de Colombia viven también decenas de especies de abejas
nativas que no tienen un aguijón capaz de picar, y que los campesinos conocen con nombres como angelitas o abejas
monas.

\pa Estas abejas son, en general, más pequeñas y producen menos miel que las comunes. Su miel, sin embargo, es muy
apreciada: es más líquida, algo ácida y, en muchas regiones, se usa como remedio casero. Como no pican, pueden criarse
cerca de las casas sin peligro, lo que ha despertado el interés de familias que buscan un ingreso adicional.

\pa Su importancia va más allá de la miel. Muchas plantas nativas dependen de ellas para la polinización, porque sus
flores tienen formas o tamaños que otras abejas aprovechan mal. Cuando un bosque pierde a sus abejas nativas, algunas de
esas plantas producen menos frutos y semillas.

\pa El principal riesgo para ellas no es una enfermedad sino la pérdida de lugares donde anidar: viven en troncos huecos y
en grietas de árboles viejos, que desaparecen cuando se tala el bosque. Por eso, quienes las crían en cajas ayudan, sin
proponérselo, a conservar poblaciones que de otro modo podrían reducirse.
\end{texto}

\subsection{Preguntas comentadas}
\begin{comentada}{¿Cuál es una paráfrasis correcta de la primera oración del \S\,2?}
\mal{A}{Todas las abejas nativas producen más miel que las comunes.}{\emph{Dice lo contrario}: producen menos.}
\bien{B}{Por lo general, son de menor tamaño y dan menos miel que las abejas comunes.}{Conserva «en general», «más
pequeñas» y «menos miel».}
\mal{C}{Son siempre más pequeñas y nunca producen miel.}{\emph{Dice de más}: cambia «en general» por «siempre» y «menos
miel» por «nunca».}
\mal{D}{Las abejas comunes son más pequeñas que las nativas.}{Intercambia los papeles.}
\end{comentada}

\taller
\piso{Calienta}{reescribir}
\begin{ejercicios}
\item Parafrasea: «El principal riesgo para ellas no es una enfermedad sino la pérdida de lugares donde anidar». \sol{Por ejemplo: «Lo que más las amenaza es quedarse sin sitios para hacer sus nidos, más que las enfermedades».}
\end{ejercicios}
\piso{Como en el examen}{opción múltiple}
\begin{ejercicios}
\item ¿Cuál expresa lo mismo que «muchas plantas nativas dependen de ellas para la polinización» \pf{3}?
\opcionesl{Todas las plantas necesitan abejas nativas.}{Numerosas plantas nativas necesitan a estas abejas para polinizarse.}{Las abejas nativas dependen de las plantas.}{Algunas plantas polinizan a las abejas.}
\sol{B.}
\item «Como no pican, pueden criarse cerca de las casas sin peligro» \pf{2} equivale a:
\opcionesl{Se crían cerca de las casas porque son peligrosas.}{Su falta de aguijón permite tenerlas cerca de las viviendas sin riesgo.}{Todas las familias crían abejas cerca de sus casas.}{No pican porque viven cerca de las casas.}
\sol{B. D invierte la causa.}
\item Según el \S\,4, quienes crían abejas nativas
\opcionesl{lo hacen para conservarlas.}{contribuyen a conservarlas aunque no sea su intención.}{talan los árboles viejos.}{las enferman.}
\sol{B. «sin proponérselo».}
\end{ejercicios}

\begin{retoresuelto}
«No todas pican» significa que hay al menos algunas que no pican. \textbf{Opción B.} La A («ninguna») es la trampa clásica
del cuantificador.
\end{retoresuelto}
\bitacora

% ------------------------------------------------------------------
\leccion{pu-ie-3}{Lo que dice y lo que no}{Juego de clasificar afirmaciones: el texto lo dice, lo niega o no lo dice.}

\begin{reto}
Texto: «El estudio encontró que los estudiantes que dormían menos de seis horas obtenían notas más bajas.» ¿Cuál afirmación
\emph{no} está respaldada por el texto?
\opcionesl{Hubo un estudio sobre sueño y notas.}{Dormir poco se asoció con notas más bajas.}{Dormir poco es la causa de las malas notas.}{El estudio incluyó a estudiantes.}
\end{reto}

\tecnica{La técnica: tres cajas}
Ante cada afirmación hay tres posibilidades, no dos: el texto la \clave{dice} (o la implica sin duda), la \clave{niega}, o
\clave{no se pronuncia}. La tercera caja es la que más se olvida: muchas opciones falsas no contradicen el texto, solo
dicen algo que el texto no afirma. Un caso frecuente es convertir una relación («los que duermen poco sacan peores
notas») en una causa («dormir poco causa malas notas»).

\figura{trescajas}{Frente a cada opción hay tres cajas, no dos.}

\begin{regla}
\begin{itemize}
  \item Clasifica: lo dice / lo niega / no lo dice.
  \item «Se asocia», «coincide», «los que… suelen…» no es lo mismo que «causa».
  \item Lo que sabes del mundo no cuenta: solo lo que el texto sostiene.
\end{itemize}
\end{regla}

\begin{texto}{Trabajar desde la sala}{Artículo · texto original escrito para este libro}
\pa Desde hace unos años, muchas empresas permiten a sus empleados trabajar algunos días desde la casa. La medida, que
empezó como una emergencia, se volvió para muchos una costumbre, y hoy divide opiniones.

\pa Quienes la defienden señalan el tiempo que se ahorra en transporte, que en ciudades grandes puede pasar de dos horas
diarias, y la posibilidad de organizar mejor la jornada. Algunas encuestas a trabajadores muestran que la mayoría prefiere
un esquema mixto: unos días en la oficina y otros en la casa.

\pa Los críticos advierten otros efectos. Sin la separación física entre el trabajo y el hogar, la jornada tiende a
alargarse; los mensajes llegan a cualquier hora. Además, los empleados nuevos aprenden más despacio cuando no pueden
observar a sus compañeros más experimentados.

\pa Hay también una desigualdad menos visible: no todos los oficios pueden hacerse a distancia, y no todos los hogares
tienen un espacio tranquilo o una buena conexión. Para quien trabaja desde la mesa del comedor con tres hijos alrededor,
la casa no es necesariamente un lugar más cómodo que la oficina.
\end{texto}

\subsection{Preguntas comentadas}
\begin{comentada}{Según el texto, ¿cuál afirmación es verdadera?}
\mal{A}{Trabajar desde la casa mejora siempre la productividad.}{\emph{No lo dice} (y además dice de más).}
\bien{B}{Según algunas encuestas, la mayoría de trabajadores prefiere combinar oficina y casa.}{\S\,2, con el matiz
«algunas encuestas».}
\mal{C}{Todos los oficios pueden hacerse a distancia.}{\emph{Dice lo contrario} \pf{4}.}
\mal{D}{Las empresas van a prohibir el trabajo en casa.}{\emph{No lo dice}.}
\end{comentada}

\taller
\piso{Calienta}{tres cajas}
\begin{ejercicios}
\item Clasifica (dice / niega / no dice): a) la jornada en casa tiende a alargarse; b) los empleados nuevos aprenden más rápido en casa; c) el trabajo en casa empezó en 2010. \sol{a) dice \pf{3}; b) niega \pf{3}; c) no dice.}
\end{ejercicios}
\piso{Como en el examen}{opción múltiple}
\begin{ejercicios}
\item Según el texto, ¿cuál es un argumento de los críticos?
\opcionesl{El transporte es muy caro.}{Los nuevos aprenden más lento sin ver a los experimentados.}{La mayoría prefiere la oficina.}{Las casas son más cómodas.}
\sol{B. \S\,3.}
\item El texto NO afirma que
\opcionesl{el trabajo en casa empezó como emergencia.}{en ciudades grandes el transporte puede pasar de dos horas diarias.}{trabajar en casa reduce los salarios.}{no todos los hogares tienen buena conexión.}
\sol{C. No se menciona.}
\end{ejercicios}

\begin{retoresuelto}
El texto informa una asociación («los que dormían menos… obtenían notas más bajas»), no una causa. Afirmar que dormir poco
\emph{causa} las malas notas va más allá de lo dicho. \textbf{Opción C.}
\end{retoresuelto}
\bitacora

% ------------------------------------------------------------------
\leccion{pu-ie-4}{Quién dice qué}{Juego de asignar cada cita a su voz.}

\begin{reto}
En un artículo, el periodista escribe: «Según el ingeniero, la ciclorruta reducirá los trancones; los comerciantes, en
cambio, temen perder clientes». ¿Quién teme perder clientes?
\opcionesd{el periodista}{el ingeniero}{los comerciantes}{los ciclistas}
\end{reto}

\tecnica{La técnica: un color por voz}
Muchos textos reúnen varias \clave{voces}: el autor, los expertos que cita, los testigos, los opositores. Las preguntas
del examen preguntan a menudo qué piensa \emph{cada uno}, y los distractores atribuyen a uno lo que dijo otro. Lo práctico
es marcar al margen, en cada párrafo, quién habla: A (autor), E (experto), C (comerciante)…

Las marcas de una voz ajena son: comillas, verbos de decir («afirma», «advierte», «sostiene»), expresiones como «según»,
«para», «en opinión de».

\figura{voces}{Cada idea tiene dueño: un color por voz.}

\begin{regla}
\begin{itemize}
  \item Marca al margen quién habla en cada parte.
  \item «Según X» y las comillas: la idea es de X, no necesariamente del autor.
  \item Antes de elegir, verifica que la opción atribuya la idea a la persona correcta.
\end{itemize}
\end{regla}

\begin{texto}{La ciclorruta de la calle Bolívar}{Noticia con voces · texto original escrito para este libro (personas imaginarias)}
\pa La alcaldía anunció que en marzo empezará la construcción de una ciclorruta de dos kilómetros en la calle Bolívar, uno
de los ejes comerciales del municipio. El proyecto eliminará uno de los dos carriles para carros.

\pa Para la ingeniera de movilidad Sandra Ríos, la obra «no le quita espacio a nadie: se lo devuelve a la mayoría». Según
sus conteos, en horas pico pasan por la calle más personas a pie y en bicicleta que en carro.

\pa Los comerciantes no están convencidos. «Nuestros clientes vienen en carro y necesitan parquear», dijo Julio Cárdenas,
dueño de una ferretería. Teme que las ventas caigan durante la obra y después de ella.

\pa Laura Mesa, del colectivo Bici Bolívar, recordó que en otras ciudades los negocios cercanos a las ciclorrutas «al
principio protestaron y después pidieron que las extendieran». Reconoció, sin embargo, que la obra debe hacerse rápido
para no ahogar al comercio.

\pa La alcaldía prometió que la construcción durará como máximo tres meses y que habrá zonas de cargue y descargue cada
doscientos metros.
\end{texto}

\subsection{Preguntas comentadas}
\begin{comentada}{¿Quién reconoce que la obra debe hacerse rápido para no perjudicar al comercio?}
\mal{A}{Julio Cárdenas.}{Él teme la caída de ventas, pero esa frase no es suya.}
\bien{B}{Laura Mesa.}{\S\,4: «Reconoció, sin embargo, que la obra debe hacerse rápido».}
\mal{C}{Sandra Ríos.}{Ella habla de quién usa la calle \pf{2}.}
\mal{D}{La alcaldía.}{La alcaldía promete un plazo \pf{5}, que no es lo mismo.}
\end{comentada}

\taller
\piso{Como en el examen}{opción múltiple}
\begin{ejercicios}
\item Según la ingeniera, en horas pico por la calle Bolívar pasan
\opcionesl{más carros que personas.}{más personas a pie y en bicicleta que en carro.}{solo ciclistas.}{pocos clientes.}
\sol{B. \S\,2.}
\item La idea de que los negocios «después pidieron que las extendieran» \pf{4} se refiere a
\opcionesl{los comerciantes de la calle Bolívar.}{negocios de otras ciudades, según Laura Mesa.}{la opinión del periodista.}{la alcaldía.}
\sol{B.}
\item ¿Cuál de estas afirmaciones es del periodista y no de una persona citada?
\opcionesl{«no le quita espacio a nadie»}{El proyecto eliminará uno de los dos carriles para carros.}{«Nuestros clientes vienen en carro»}{Las ventas caerán.}
\sol{B. \S\,1, sin comillas ni verbo de decir.}
\end{ejercicios}

\begin{retoresuelto}
La frase dice «los comerciantes, en cambio, temen perder clientes». \textbf{Opción C.} El periodista solo reporta.
\end{retoresuelto}
\bitacora

% ------------------------------------------------------------------
\evaluacion
\begin{texto}{El mercado que despierta a las tres}{Crónica · texto original escrito para este libro}
\pa A las tres de la mañana, cuando la ciudad todavía duerme, la plaza mayorista ya está despierta. Los primeros camiones
llegan del oriente cargados de papa, zanahoria y cebolla; un poco después entran los del suroeste, con plátano y café.

\pa Hernán, cotero desde hace veinte años, carga en la espalda bultos de cincuenta kilos. Dice que el trabajo no es lo
más duro: lo más duro es que los precios cambian cada día y nadie sabe por qué. «Un lunes la papa vale una cosa y el
martes otra, y el campesino siempre pierde», asegura.

\pa Doña Marleny, que vende frutas en un puesto desde hace treinta años, no está de acuerdo del todo. Para ella, los precios
dependen sobre todo del clima: «Cuando llueve mucho en el oriente, llega menos papa y sube». Admite, sin embargo, que los
intermediarios se quedan con una parte grande de la ganancia.

\pa Hacia las siete, los compradores de los supermercados ya se han ido y llegan las señoras de los barrios, que buscan
lo que sobra a mejor precio. A las diez, la plaza empieza a vaciarse. Hernán se sienta en un bulto vacío, se toma un
tinto y cuenta los billetes de la jornada.
\end{texto}
\begin{ejercicios}
\item Según el texto, ¿de dónde llegan los camiones con plátano y café?
\opcionesd{del oriente}{del suroeste}{de los barrios}{de los supermercados}
\sol{B. \S\,1. (3.1)}
\item ¿Quién afirma que el campesino siempre pierde?
\opcionesd{doña Marleny}{Hernán}{el cronista}{los intermediarios}
\sol{B. \S\,2. (3.4)}
\item Según doña Marleny, los precios dependen sobre todo
\opcionesd{de los intermediarios}{del clima}{de los supermercados}{de los coteros}
\sol{B. Aunque admite que los intermediarios ganan mucho. (3.4)}
\item «Doña Marleny no está de acuerdo del todo» \pf{3} significa que
\opcionesl{rechaza por completo lo que dice Hernán.}{coincide en parte con Hernán, pero explica los precios de otra forma.}{nunca habla con Hernán.}{está de acuerdo en todo.}
\sol{B. (3.2)}
\item El texto NO dice que
\opcionesl{la plaza empieza a vaciarse a las diez.}{las señoras de los barrios buscan lo que sobra a mejor precio.}{los coteros ganan más que los vendedores.}{Hernán carga bultos de cincuenta kilos.}
\sol{C. (3.3)}
\end{ejercicios}
